\section*{Bab \ref{ch:g12:cells-and-electrolysis} --- Sel Elektrokimia dan Elektrolisis}

\begin{solution}{exo:g12:cells-and-electrolysis:1}
Elektrode seng adalah \omterm{def:g12:cells-and-electrolysis:anode}{anode}, \ce{Zn -> Zn^{2+} + 2e-}; elektrode
tembaga adalah \omterm{def:g12:cells-and-electrolysis:anode}{katode}, \ce{Cu^{2+} + 2e- -> Cu}.
\end{solution}

\begin{solution}{exo:g12:cells-and-electrolysis:2}
% ledger: const:F
$Q = 0.20 \times 1800 = \qty{360}{C}$; $n(e^-) = 360/96485 =
\qty{3.7e-3}{mol}$.
\end{solution}

\begin{solution}{exo:g12:cells-and-electrolysis:3}
\omterm{def:g12:cells-and-electrolysis:cell}{Jembatan garam} menutup rangkaian dan menjaga setiap larutan tetap netral
secara listrik, dengan \omterm{def:g9:ions:ion}{ion-ionnya} bergerak ke dalam ruang-ruang itu.
Tanpanya, rangkaiannya terbuka: tidak ada arus yang mengalir,
voltmeter menunjukkan nol, dan lampu tetap padam.
\end{solution}

\begin{solution}{exo:g12:cells-and-electrolysis:4}
$2.5 \times 3600 = \qty{9000}{C}$.
\end{solution}

\begin{solution}{exo:g12:cells-and-electrolysis:5}
Tidak: reaksinya berlangsung ke arah yang berlawanan dengan arah
spontannya. Generator memasok energinya.
\end{solution}

\begin{solution}{exo:g12:cells-and-electrolysis:6}
\omterm{def:g12:cells-and-electrolysis:anode}{Anode} (tembaga, $-$): \ce{Cu -> Cu^{2+} + 2e-}. \omterm{def:g12:cells-and-electrolysis:anode}{Katode} (perak, $+$):
\ce{Ag+ + e- -> Ag}. Keseluruhan: \ce{Cu + 2Ag+ -> Cu^{2+} + 2Ag}.
\end{solution}

\begin{solution}{exo:g12:cells-and-electrolysis:7}
Tiga tahun sama dengan $3 \times 365 \times 24 = \qty{26280}{h}$;
$0.010 \times 26280 = \qty{263}{mA.h}$.
\end{solution}

\begin{solution}{exo:g12:cells-and-electrolysis:8}
% ledger: const:F, aw:Ag
$Q = 0.50 \times 2400 = \qty{1200}{C}$; $n(e^-) = n(\ce{Ag}) =
1200/96485 = \qty{0.0124}{mol}$; $m = 0.0124 \times 107.9 =
\qty{1.34}{g}$.
\end{solution}

\begin{solution}{exo:g12:cells-and-electrolysis:9}
Arus di dalam kawat mengalir dari tembaga ($+$) ke seng ($-$),
berlawanan dengan \omterm{def:g9:inside-the-atom:nucleus}{elektron}. \omterm{def:g9:ions:ion}{Ion} \ce{K+} bergerak menuju ruang tembaga,
yang kehilangan \ce{Cu^{2+}}; \omterm{def:g9:ions:ion}{ion} \ce{NO3-} menuju ruang seng, yang
mendapat tambahan \ce{Zn^{2+}}.
\end{solution}

\begin{solution}{exo:g12:cells-and-electrolysis:10}
Pada keduanya, \omterm{def:g12:cells-and-electrolysis:anode}{anode} adalah tempat terjadinya \omterm{def:g11:redox:oxidation}{oksidasi}. Pada sel, \omterm{def:g12:cells-and-electrolysis:anode}{anode}
bertanda $-$, reaksinya spontan, dan energi kimia berubah menjadi energi
listrik; pada \omterm{def:g12:cells-and-electrolysis:electrolysis}{elektrolisis}, \omterm{def:g12:cells-and-electrolysis:anode}{anode} dihubungkan ke kutub $+$ generator,
reaksinya dipaksakan, dan energi listrik berubah menjadi energi kimia.
\end{solution}

\begin{solution}{exo:g12:cells-and-electrolysis:11}
\qty{12}{mL}: persamaan \ce{2H2O -> 2H2 + O2} memberikan dua \omterm{def:g7:atoms-and-molecules:molecule}{molekul}
hidrogen untuk satu \omterm{def:g7:atoms-and-molecules:molecule}{molekul} oksigen, dan jumlah gas yang sama menempati
volume yang sama.
\end{solution}

\begin{solution}{exo:g12:cells-and-electrolysis:12}
% ledger: aw:Zn, const:F
$n(\ce{Zn}) = 5.0/65.4 = \qty{0.076}{mol}$;
$n(\ce{Cu^{2+}}) = 0.100 \times 0.50 = \qty{0.050}{mol}$; keduanya
bereaksi satu banding satu, jadi \omterm{def:g9:ions:ion}{ion-ion} tembaga habis lebih dahulu.
$n(e^-) = 2 \times 0.050
= \qty{0.100}{mol}$, $Q = \qty{9.65e3}{C} = \qty{2.68}{A.h}$. Pada
\qty{50}{mA}: $2.68/0.050 = \qty{54}{h}$.
\end{solution}

\begin{solution}{exo:g12:cells-and-electrolysis:13}
$Q = 2.0 \times 3600 = \qty{7200}{C}$; $E = 7200 \times 1.5 =
\qty{1.08e4}{J}$, yaitu $1.5 \times 2.0 = \qty{3.0}{W.h}$.
\end{solution}

\begin{solution}{exo:g12:cells-and-electrolysis:14}
% ledger: const:F
$Q = \qty{7200}{C}$, $n(e^-) = \qty{0.0746}{mol}$. Dua \omterm{def:g9:inside-the-atom:nucleus}{elektron} per
\ce{H2}: $n(\ce{H2}) = \qty{0.0373}{mol}$, $V = \qty{0.90}{L}$. Empat
per \ce{O2}: $n(\ce{O2}) = \qty{0.0187}{mol}$, $V = \qty{0.45}{L}$.
\end{solution}

\begin{solution}{exo:g12:cells-and-electrolysis:15}
Tembaga yang \omterm{def:g2:mixing-and-dissolving:dissolve}{larut} di \omterm{def:g12:cells-and-electrolysis:anode}{anode} sama dengan tembaga yang diendapkan, karena
\omterm{def:g9:inside-the-atom:nucleus}{elektron} yang sama melewati keduanya: $2.84/3.00 = \qty{94.7}{\percent}$
tembaga.
\end{solution}

\begin{solution}{pb:g12:cells-and-electrolysis:1}
% ledger: const:F, const:NA, aw:Li, dch:CH4
\textbf{1.} \omterm{def:g12:cells-and-electrolysis:anode}{Anode}: litium teroksidasi di sana.

\textbf{2.} Elektrode grafit, yaitu \omterm{def:g12:cells-and-electrolysis:anode}{anode}.

\textbf{3.} \omterm{def:g9:inside-the-atom:nucleus}{Elektron} berjalan melalui ponsel dari elektrode grafit ke
elektrode oksida; \omterm{def:g9:ions:ion}{ion-ion} litium menyeberangi elektrolit ke arah yang
sama, dari grafit ke oksida.

\textbf{4.} Reaksinya dapat dipaksa berbalik arah oleh pengisi daya.

\textbf{5.} \qty{4.000}{A.h}, yaitu $4.000 \times 3600 =
\qty{14400}{C}$.

\textbf{6.} $14400/96485 = \qty{0.1492}{mol}$.

\textbf{7.} $0.1492 \times \num{6.022e23} = \num{8.99e22}$ \omterm{def:g9:inside-the-atom:nucleus}{elektron}.

\textbf{8.} $4.000/0.25 = \qty{16}{h}$.

\textbf{9.} $0.20 \times 14400 = \qty{2880}{C}$.

\textbf{10.} $E = 14400 \times 3.7 = \qty{5.3e4}{J}$, sekitar
\qty{53}{kJ}.

\textbf{11.} $53280/3600 = \qty{14.8}{W.h}$.

\textbf{12.} \omterm{def:g4:what-a-fire-needs:burning}{Pembakaran} \qty{1}{g} metana melepaskan energi sedikit
lebih banyak (\qty{55.7}{kJ}) daripada yang disimpan baterai ponsel yang
penuh.

\textbf{13.} $1000/14.8 = 67.6$: 67 kali pengisian penuh.

\textbf{14.} \ce{Li+ + e- -> Li}: sebuah \omterm{def:g11:redox:oxidation}{reduksi}.

\textbf{15.} Pengisi daya memaksa reaksinya ke arah yang berlawanan
dengan arah spontannya.

\textbf{16.} $14400/2.0 = \qty{7200}{s}$, \qty{2.0}{h}.

\textbf{17.} Satu \omterm{def:g9:ions:ion}{ion} litium per \omterm{def:g9:inside-the-atom:nucleus}{elektron}: \qty{0.1492}{mol}.

\textbf{18.} $0.1492 \times 6.9 = \qty{1.03}{g}$.

\textbf{19.} Sekitar \qty{1.03}{g} litium bolak-balik di antara
elektrode-elektrode dalam satu pengisian penuh.
\end{solution}
